\chapter{Fase 7 -- Escrita dos arquivos e execução do modelo}
\label{cap:fase7}

\section{Composição do \texttt{init.nc} completo}

O \texttt{init.nc} real do \mpasa{} contém 135 variáveis. Das fases
anteriores (1 a 6), o conjunto de programas deste trabalho calcula e
escreve diretamente 54 delas -- a grade vertical (Fase~2), os campos de
massa e vento (Fases~3--4), e os campos de superfície/solo (Fase~6). As
81 variáveis restantes são de \textbf{cópia direta} do arquivo de malha
estática (\texttt{static.nc}): geometria e conectividade da malha
(coordenadas de célula/aresta/vértice, tabelas de vizinhança), uso do
solo, orografia bruta, e coeficientes numéricos fixos da malha (não
dependentes do dado meteorológico nem do tempo de execução). Essa
composição é feita por mesclagem de arquivos NetCDF
(\texttt{ncks~-A}, utilitário do pacote NCO), não por um terceiro
programa Fortran -- uma escolha de engenharia que evita duplicar, em
Fortran, a leitura e reescrita de 81 campos que já existem prontos, sem
nenhuma transformação necessária.

\begin{figure}[H]
\centering
\begin{tikzpicture}[
    node distance=8mm and 12mm,
    box/.style={draw, rounded corners, minimum height=0.9cm, minimum width=3.2cm, align=center, font=\footnotesize, fill=blue!6},
    arr/.style={-{Stealth[length=2mm]}, thick},
  ]
  \node[box] (static) {\texttt{<REGIÃO>.static.nc}\\(81 campos estáticos)};
  \node[box, right=of static, fill=green!10] (native) {\texttt{gen\_init\_native}\\(Fases 2+3+4+6)};
  \node[box, below=of native, fill=orange!10] (computed) {\texttt{computed\_fields.nc}\\(54 campos calculados)};
  \node[box, right=of computed, fill=red!8] (merge) {\texttt{ncks -A}\\(mesclagem)};
  \node[box, right=of merge] (final) {\texttt{init.nc}\\(135 campos)};

  \draw[arr] (static) -- (native);
  \draw[arr] (native) -- (computed);
  \draw[arr] (static.south) |- ($(merge.west)+(0,0.35)$);
  \draw[arr] (computed) -- (merge);
  \draw[arr] (merge) -- (final);
\end{tikzpicture}
\caption{Composição do \texttt{init.nc} final: os 54 campos calculados
  pelos programas desta rota são mesclados, via \texttt{ncks -A}, com
  os 81 campos de cópia direta do arquivo de malha estática.}
\label{fig:composicao_init}
\end{figure}

O \texttt{lbc.*.nc} tem um esquema mais simples: apenas os sete campos
prognósticos (\texttt{lbc\_qv}, \texttt{lbc\_qc}, \texttt{lbc\_qr},
\texttt{lbc\_u}, \texttt{lbc\_w}, \texttt{lbc\_rho}, \texttt{lbc\_theta}),
um arquivo por tempo de fronteira, sem geometria nenhuma -- a malha e a
grade vertical são lidas pelo \texttt{mpas\_atmosphere} do próprio
\texttt{init.nc}, via o \emph{stream} de entrada. O programa
\texttt{gen\_lbc\_native} reaproveita, por isso, a grade vertical já
calculada e gravada no \texttt{init.nc} (não a recalcula), rodando apenas
as Fases~3--4 para cada tempo requisitado.

\section{Dois bugs de formato encontrados apenas na execução real}
\label{sec:bugs_formato}

A validação numérica campo a campo contra arquivos de referência
(\cref{cap:resultados}), por mais rigorosa que tenha sido, não é
suficiente para garantir que um arquivo NetCDF seja de fato
\emph{consumível} pelo programa de destino -- um ponto que se revelou
concretamente ao tentar, pela primeira vez, alimentar o
\texttt{mpas\_atmosphere} real com os arquivos gerados por esta rota.

\begin{caixaachado}
\textbf{Ausência da variável \texttt{xtime}.} A primeira tentativa de
execução falhou imediatamente, antes mesmo do primeiro passo de tempo,
com o erro \texttt{"ERROR: File SouthAmerica.init.nc does not contain
the time 2026-01-01\_00:00:00"}. O \texttt{init.nc} gerado continha a
variável \texttt{initial\_time} (uma string informativa, escrita
corretamente), mas não a variável \texttt{xtime(Time, StrLen)} -- que é,
na verdade, a que o \texttt{mpas\_atmosphere} lê para confirmar que o
arquivo de entrada contém o tempo solicitado pelo
\texttt{config\_start\_time}. A variável havia sido listada, na fase de
levantamento inicial das 135 variáveis do \texttt{init.nc} real, mas não
percebida como funcionalmente distinta de \texttt{initial\_time} até
esse ponto -- um lembrete de que um levantamento de nomes de variável não
substitui saber \emph{qual delas o consumidor realmente lê}.

\textbf{Lixo de memória no \texttt{xtime} do \texttt{lbc.*.nc}.} Corrigido
o problema acima, uma segunda falha ocorreu já dentro do primeiro passo
de tempo, com um erro de execução do próprio \emph{runtime} Fortran
(\texttt{"Fortran runtime error: Bad integer for item 1 in list input"})
dentro da rotina de gerenciamento de tempo do framework
(\texttt{mpas\_timekeeping.F}), derrubando simultaneamente todos os 32
processos MPI. A causa, encontrada inspecionando os \emph{bytes} brutos
do arquivo gerado, foi uma chamada \texttt{nf90\_put\_var} recebendo uma
expressão de string de 24 caracteres para um campo NetCDF declarado com
64 bytes: os 40 bytes restantes ficavam preenchidos com lixo de memória
não inicializado (confirmado como binário arbitrário, não espaços em
branco), que o \emph{parser} de data do \textit{framework} -- uma leitura
formatada por posição fixa de coluna, rígida por natureza -- não
conseguia interpretar. A correção foi declarar explicitamente uma
variável Fortran do tipo \texttt{character(len=64)} antes de passá-la ao
\texttt{nf90\_put\_var} (o Fortran preenche automaticamente o excedente
com espaços em branco ao atribuir uma string mais curta a uma variável de
tamanho fixo maior), em vez de passar a expressão concatenada
diretamente.
\end{caixaachado}

Um terceiro problema, de natureza puramente operacional (não um bug do
código Fortran), também surgiu apenas na execução real: o script de
orquestração da rodada carrega um arquivo de ambiente do cluster que
executa \texttt{conda deactivate} internamente; quando o script é
invocado como \texttt{bash script.bash} (\emph{shell} não interativo), o
\texttt{\textasciitilde/.bashrc} não é lido por padrão mesmo com o
\textit{hook} do \texttt{conda} já inicializado uma vez no terminal
interativo do usuário, e o comando falha, derrubando o script inteiro
(\texttt{set -e}). A correção carregou defensivamente o \textit{hook} do
\texttt{conda} (\texttt{etc/profile.d/conda.sh}) no início do script, se
necessário, antes de delegar ao arquivo de ambiente do cluster.

A lição transversal aos três problemas: nenhum deles era de física ou de
numérica -- o \texttt{init.nc}/\texttt{lbc.*.nc} já batiam campo a campo
com a referência antes de qualquer um deles ser encontrado -- e mesmo
assim o modelo real recusava ou travava com os arquivos. Validação
numérica contra referência não substitui validar contra o consumidor
real; um arquivo pode estar fisicamente correto e ainda assim ser
inutilizável por um detalhe de formato invisível a uma comparação
campo a campo.

\section{Execução completa}

Corrigidos os três problemas, a execução do \texttt{mpas\_atmosphere}
real, com 32 processos MPI, a partir do \texttt{init.nc}/\texttt{lbc.*.nc}
gerados inteiramente por esta rota, completou 240 passos de tempo (24
horas de integração, malha \texttt{SouthAmerica}, $\sim\SI{60}{\kilo\metre}$
de resolução) em aproximadamente 106 segundos de tempo de parede, sem
nenhuma mensagem de erro ou erro crítico (apenas 24 avisos inofensivos,
relativos a variáveis de gelo/graupel/neve ausentes do
\textit{first-guess} -- esperado, já que o dado de origem não modela
esses hidrometeoros). A condição de contorno lateral foi aplicada
corretamente nos quatro intervalos de seis horas cobertos pelos arquivos
\texttt{lbc.*.nc}, gerando 25 arquivos de saída
(\texttt{history.*.nc}, um por hora). Os resultados quantitativos dessa
execução, e sua comparação com a rodada de referência gerada pela rota
original, são apresentados no \cref{cap:resultados}.
